\chapter{Riscos}

Conforme a Resolução CNS 466/12, não existem pesquisas sem riscos aos participantes. Os riscos envolvem todas as áreas inerentes à vida humana, incluindo danos às dimensões física, psíquica, moral, intelectual, social e profissional. A seguir, os riscos identificados e as estratégias de minimização.

\begin{table}[htbp]
    \centering
    \caption{\label{tab:riscos}Riscos identificados e formas de minimização}
    \begin{tabular}{|p{3.2cm}|p{5.5cm}|p{6cm}|}
    \hline
    \textbf{Risco} & \textbf{Descrição} & \textbf{Forma de Minimização} \\ \hline
    Possível frustração ou desconforto ao errar questões durante o uso da plataforma. & O sistema de gamificação é projetado para tornar o erro parte do aprendizado (penalidades leves de \(-2\%\)), com dicas progressivas que incentivam a persistência sem punição excessiva. \\ \hline
    Constrangimento & Possível constrangimento ao responder formulários de avaliação sobre desempenho. & Formulários anônimos; Escala Likert de 5 pontos para respostas rápidas; tempo de preenchimento < 1 minuto. \\ \hline
    Cansaço & Possível fadiga cognitiva durante sessões prolongadas de uso da plataforma. & Sessões de uso com duração controlada; interface gamificada que mantém o engajamento; interrupção voluntária a qualquer momento. \\ \hline
    Risco à privacidade & Possível exposição de dados pessoais durante a coleta de dados. & Dados coletados de forma anônima; ausência de coleta de dados identificáveis; armazenamento seguro em banco de dados; observância da Lei Geral de Proteção de Dados (Lei 13.709/18). \\ \hline
    Risco moral & Possível uso indevido dos dados da pesquisa em detrimento dos participantes. & Compromisso ético do pesquisador de utilizar os dados exclusivamente para fins acadêmicos; confidencialidade garantida; dados descartados após a conclusão da pesquisa. \\ \hline
    \end{tabular}
    \fonte{Elaborado pelo autor (2026).}
\end{table}

\textbf{Ponderação entre riscos e benefícios:} conforme orientação do CEP/UNEB, se os riscos excederem os benefícios, não se justifica a execução da pesquisa. Neste projeto, os riscos são mínimos e todos passíveis de minimização, enquanto os benefícios são relevantes para os participantes e para a sociedade, justificando a execução da pesquisa.